\documentclass[11pt]{article}
\usepackage[margin=1in]{geometry}
\usepackage{amsmath,amssymb,amsthm,graphicx}
\usepackage[hidelinks]{hyperref}
\setlength{\parindent}{0pt}
\setlength{\parskip}{0.65em}
\newtheorem*{theorem*}{Theorem}
\newtheorem*{lemma*}{Lemma}
\newtheorem*{corollary*}{Corollary}
\newcommand{\id}{\mathrm{id}}
\newcommand{\cyc}[1]{\#(#1)}
\newcommand{\dist}{d}

\begin{document}
\begin{center}
{\Large Partial Result Packet}\\[0.2em]
{\large A two-cycle obstruction to bounded geodesic repair}\\[0.3em]
\textbf{Status:} partial result, likely valid
\end{center}

\section*{Source problem and scope}

The source is C\'ecilia Lancien,
\emph{High dimension and symmetries in quantum information theory},
arXiv:1607.06843. A local copy is \texttt{source\_paper.pdf}.
For permutations on \(p\) points, write
\[
 \dist(\sigma,\tau)=p-\cyc{\sigma^{-1}\tau}
\]
for the transposition Cayley distance. The source asks whether excess
\(2\delta\) over a geodesic can be repaired within distance \(O(\delta)\),
with polynomially bounded fibres.

\begin{center}
\includegraphics[width=0.92\linewidth]{figures/open_problem_crop.png}
\end{center}

The displayed question is for one canonical full cycle \(\gamma\). The
preceding paragraph says that the case truly needed in the application is
\(\gamma_f\), the product of canonical cycles on the level sets of \(f\).
This packet disproves that multi-cycle extension with two level sets.  A
sibling counterexample packet produced later in the run settles the printed
single-full-cycle question negatively.

\clearpage
\section*{Proof intuition}

Take two cyclic blocks \(A\) and \(B\). Let \(\gamma_f\) rotate each block,
while \(\alpha\) pairs points across the blocks in reverse cyclic order.
Both \(\alpha\) and \(\alpha^{-1}\gamma_f\) are perfect matchings, so the
defect is one. However, every permutation on a shortest path from \(\id\)
to \(\gamma_f\) preserves both blocks, while \(\alpha\) swaps them. The
quotient must alternate between the blocks, giving linear distance.

\section*{Interval lemma}

\begin{lemma*}
If \(\dist(\id,\pi)+\dist(\pi,\gamma)=\dist(\id,\gamma)\), then every cycle
of \(\pi\) is contained in a cycle of \(\gamma\).
\end{lemma*}
\begin{proof}
Concatenate shortest paths from \(\id\) to \(\pi\) and from \(\pi\) to
\(\gamma\). The equality makes this a shortest path. At every step from the
identity, the transposition must increase length by one, hence merge two
cycles. A split would decrease length. Thus cycle partitions only become
coarser, proving the claim.
\end{proof}

\section*{Main result}

\begin{theorem*}
For every \(n\geq2\), there are \(\gamma,\alpha\in S_{2n}\) such that
\(\gamma\) is a product of two \(n\)-cycles,
\[
 \dist(\id,\alpha)+\dist(\alpha,\gamma)
 =\dist(\id,\gamma)+2,
\]
and
\[
 \dist\bigl(\alpha,[\id,\gamma]\bigr)=n.
\]
\end{theorem*}

\begin{proof}
Use subscripts in \(\mathbb Z/n\mathbb Z\). On
\[
 A=\{a_i:i\in\mathbb Z/n\mathbb Z\},\qquad
 B=\{b_i:i\in\mathbb Z/n\mathbb Z\},
\]
define
\[
 \gamma=(a_0\ a_1\ \cdots\ a_{n-1})
        (b_0\ b_1\ \cdots\ b_{n-1}),\qquad
 \alpha=\prod_i(a_i\ b_{-i}).
\]
The permutation \(\alpha\) is a fixed-point-free involution, so
\(\cyc{\alpha}=n\). Moreover
\[
 \alpha\gamma(a_i)=b_{-i-1},\qquad
 \alpha\gamma(b_j)=a_{-j-1}.
\]
Hence \(\alpha^{-1}\gamma=\alpha\gamma\) is another fixed-point-free
involution and \(\cyc{\alpha^{-1}\gamma}=n\). Therefore
\[
 \dist(\id,\alpha)=n,\qquad
 \dist(\alpha,\gamma)=n,\qquad
 \dist(\id,\gamma)=2n-2,
\]
which proves the excess-two identity.

If \(\pi\in[\id,\gamma]\), the interval lemma says that \(\pi\) preserves
\(A\) and \(B\). Since \(\alpha^{-1}=\alpha\) swaps them,
\(\alpha^{-1}\pi\) also swaps them. Every cycle of this quotient alternates
between \(A\) and \(B\), so it has at most \(n\) cycles. Consequently
\[
 \dist(\alpha,\pi)=2n-\cyc{\alpha^{-1}\pi}\geq n.
\]
Equality holds at \(\pi=\id\), proving the claimed exact distance.
\end{proof}

\begin{corollary*}
No universal \(\theta<\infty\) gives a repair bound
\(\dist(\alpha,\alpha')\leq2\theta\delta\) for every product
\(\gamma_f\) of canonical cycles.
\end{corollary*}
\begin{proof}
The source's multi-cycle defect equation is
\[
 \cyc{\gamma^{-1}\alpha}+\cyc{\alpha}
 =p+\cyc{\gamma}-2\delta.
\]
Here the left side is \(2n\), while \(p=2n\) and \(\cyc{\gamma}=2\), so
\(\delta=1\). The required distance is nevertheless at least \(n\), which is
unbounded. The additional fibre condition therefore cannot rescue the claim.
\end{proof}

\section*{Verification, limitations, and novelty}

The proof is uniform and computationally independent. The script
\texttt{code/verify\_obstruction.py} exhaustively enumerates \(S_{2n}\) for
\(2\leq n\leq4\), selects the interval by its distance equality, and returns:
\begin{verbatim}
n=2: interval_size=4, excess=2, min_distance=2
n=3: interval_size=25, excess=2, min_distance=3
n=4: interval_size=196, excess=2, min_distance=4
\end{verbatim}

The result in this packet is partial because its block obstruction uses two
cycles.  Subsequent work in the same run found a different defect-one family
for a single full cycle whose distance from the geodesic interval is
unbounded; see the sibling full-cycle counterexample packet in this run.
Thus the exact printed question is now also answered negatively.

A bounded search of the run indexes and arXiv for the source identifier,
permutation genus, absolute order, geodesic-interval distance, and two-cycle
products found related work on noncrossing and plane permutations, but no
occurrence of this explicit defect-one obstruction. Standard genus-slicing
operations may move the face cycle and do not supply the fixed-face statement.

\textbf{Human review focus.} Check the source convention
\(\cyc{\gamma_f^{-1}\alpha}+\cyc{\alpha}
=p+\cyc{\gamma_f}-2\delta\) and the interval lemma. These two points make
the scope and the defect computation transparent.

\section*{References}
\begin{itemize}
\item C. Lancien,\\
\emph{High dimension and symmetries in quantum information theory},
arXiv:1607.06843, Section 9.16.3.
\item R. X. F. Chen and C. M. Reidys, \emph{Plane permutations and
applications to a result of Zagier--Stanley and distances of permutations},
arXiv:1502.07674.
\item R. Cori and G. Hetyei, \emph{Counting genus one partitions and
permutations}, arXiv:1306.4628.
\item A. Hock, \emph{Genus permutations and genus partitions},
arXiv:2306.16237.
\end{itemize}
\end{document}
